\documentclass[final]{article}

\usepackage{APS360}
\usepackage[utf8]{inputenc} % allow utf-8 input
\usepackage[T1]{fontenc}    % use 8-bit T1 fonts
\usepackage{hyperref}       % hyperlinks
\usepackage{url}            % simple URL typesetting
\usepackage{booktabs}       % professional-quality tables
\usepackage{amsfonts}       % blackboard math symbols
\usepackage{nicefrac}       % compact symbols for 1/2, etc.
\usepackage{microtype}      % microtypography
\usepackage{xcolor}         % colors
\usepackage{graphicx}

% APS360.sty is a NeurIPS style; with [final] its first-page footer calls an
% undefined \@trackname. Blank the footer instead of editing the course file.
\makeatletter
\renewcommand{\@noticestring}{}
\makeatother

\title{Phone-Level Mispronunciation Detection for Korean Learners of English}

\author{%
  Jinwoo Lee \\
  Student Number: TODO, lee04jinwoo@gmail.com\\
  \url{https://github.com/lee04jinwoo-sys/aps360-mispronunciation}\\
}

\begin{document}

\maketitle

\vspace{-0.5in}

% Page budget: main text <= 2 pages (references excluded). 3 pages = -20%.

\section*{Introduction}
% (4 pts) Goal, why it matters, why deep learning.
% - Goal: given a recording of a read English sentence and its canonical
%   phone sequence, flag each phone as correctly pronounced or mispronounced.
% - Why: learners (esp. Korean L1, e.g. /r/-/l/, /f/-/p/, /v/-/b/, /z/-/dZ/)
%   rarely get phone-level feedback outside a classroom; human raters are scarce.
% - Why DL: acoustic realisations vary across speakers; hand-crafted
%   thresholds (GOP) generalise poorly; CNNs learn time-frequency patterns
%   directly from spectrograms.

\section*{Illustration}
% (4 pts) Pipeline figure: waveform -> phone segmentation (annotation
% boundaries) -> log-mel patch + canonical phone embedding -> CNN -> P(error).

\begin{figure}[!h]
  \centering
  \fbox{\rule[-.5cm]{0cm}{4cm} \rule[-.5cm]{4cm}{0cm}}
  \caption{TODO: overall pipeline.}
  \label{fig:pipeline}
\end{figure}

\section*{Background \& Related Work}
% (4 pts) >= 5 prior works, one or two sentences each.
% GOP \cite{witt2000phone}; DNN acoustic model + LR \cite{hu2015improved};
% CNN-RNN-CTC end-to-end MDD \cite{leung2019cnn}; wav2vec 2.0
% \cite{baevski2020wav2vec} and its MDD fine-tuning \cite{xu2021explore,peng2021study};
% L2-ARCTIC corpus \cite{zhao2018l2arctic}.

\section*{Data Processing}
% (4 pts) Must be reproducible by a classmate. Fill numbers from
% src/data_stats.py output, not by hand.
% 1. Source: L2-ARCTIC v5.0 (CC BY-NC 4.0), 24 speakers, 6 L1s.
% 2. Use the manually annotated subset (annotation/*.TextGrid).
% 3. Parse phone tier: canonical phone, perceived phone, error type (S/D/A).
% 4. Cut each phone segment (+ context window), resample to 16 kHz,
%    compute log-mel (n_mels, hop TBD), pad/crop to fixed length.
% 5. Label: 1 if substitution/deletion, 0 if correct. Insertions handled TBD.
% 6. Split by speaker (no speaker in more than one split); Korean speakers
%    held out as a dedicated test set.
% 7. Class imbalance: report error rate; weighted loss / resampling.

\section*{Architecture}
% (2 pts) 2-D CNN over log-mel patch (conv-BN-ReLU-pool blocks) +
% learned embedding of canonical phone, concatenated before FC layers,
% sigmoid output. Rough sizes OK.

\section*{Baseline Model}
% (2 pts) Reproducible: mean/std of 13 MFCC (+deltas) over the segment +
% one-hot canonical phone -> logistic regression (scikit-learn).
% Metrics: precision / recall / F1 on the error class (standard for MDD).

\section*{Ethical Considerations}
% (2 pts) Accent vs. error (risk of treating accent as "wrong");
% only 24 speakers / 6 L1s -> bias against unseen L1s, genders, ages;
% voice is biometric data (license, no redistribution);
% over-reliance on automatic feedback.

\newpage
\bibliographystyle{plain}
\bibliography{references}

\end{document}
